\chapter{Einleitung}
\label{cha:Einleitung}



Die schnelle und effiziente Kühlung von Getränken spielt sowohl im privaten als auch im professionellen Umfeld eine wichtige Rolle. Insbesondere bei Bier ist die richtige Trinktemperatur entscheidend für Geschmack, Frischeeindruck und Genussqualität. Herkömmliche Methoden – wie das Einstellen der Flasche oder Dose in einen Kühlschrank oder das Eintauchen in kaltes Wasser – sind zwar einfach umzusetzen, jedoch häufig zeitintensiv. Ziel unseres Projekts war es daher, einen kompakten, technisch innovativen Bierkühler zu entwickeln, der die Kühlzeit deutlich reduziert und dabei mit überschaubarem technischen Aufwand realisierbar ist.

Die zentrale Idee unseres Systems basiert auf einem physikalischen Grundprinzip der Wärmeübertragung: Beim Abkühlen eines ruhenden Getränks bildet sich an der Innenseite der Flasche eine isolierende Grenzschicht aus bereits abgekühlter Flüssigkeit. Diese Schicht wirkt wie eine thermische Barriere und verlangsamt den weiteren Wärmeaustausch zwischen dem wärmeren Getränkekern und der kälteren Umgebung. Unser Bierkühler verhindert die Ausbildung dieser isolierenden Schicht, indem das Bier kontinuierlich in Rotation versetzt wird. Durch die stetige Durchmischung wird die Temperatur im Inneren ausgeglichen, wodurch die Wärme effizienter nach außen abgeführt werden kann. Es handelt sich also um den effizientesten Weg Bier in möglichst geringer Zeit zu kühlen.

Technisch wird dieses Prinzip durch das Zusammenspiel mehrerer Komponenten umgesetzt. Als Herzstück wird unser Mikrocontroller verwendet, der neben der Steuerung von der Pumpe und des Motors auch den Temperatursensor an ein Display weitergibt. Die Pumpe pumpt das Wasser vom unteren des 3D Druck Behälters durch einen Wasserschlauch nach oben, damit das Wasser von oben auf die zu kühlende Flasche tropft. Der Motor betreibt die Rotation der Flasche, indem er eine Welle, welche in das 3D Druck Gehäuse integriert ist, antreibt. Um all die Komponenten im Überblick zu haben, wird ein Display verwendet, welches die Temperatur anzeigt. Durch einen Druckknopf kann der Bierkühler angehalten beziehungsweise wieder gestartet werden.  In Kombination all dieser Komponenten entsteht so ein deutlich beschleunigter Wärmeübergang zwischen dem noch warmen Bier und dem Kühlwasser.

Besonderer Wert wurde bei der Entwicklung auf eine kompakte Bauweise, eine klare Struktur der elektronischen Steuerung sowie eine stabile mechanische Umsetzung gelegt. Das 3D-gedruckte Gehäuse ermöglicht eine passgenaue Integration aller Komponenten und bietet gleichzeitig Schutz vor Feuchtigkeit und mechanischen Einwirkungen. Die modulare Bauweise erleichtert zudem Wartung, Anpassung und Weiterentwicklung des Systems. Außerdem ist unser Ziel das Bier nicht durch zu schnelle Drehung zu stark zum schäumen zu bringen, weshalb wir auf PBM setzen, um die richtige Drehgeschwindigkeit zu bekommen. 

Diese Dokumentation soll die Konzeption, Konstruktion und Funktionsweise des Bierkühlers im Detail beschreiben. Neben den physikalisch/chemischen Grundlagen werden die mechanische Konstruktion, die elektronische Steuerung, die Programmierung des Mikrocontrollers sowie mögliche Optimierungsansätze erläutert. Um in diesen vielen verschiedenen Projektteilen gut aufgestellt zu sein, haben wir für die verschiedenen Punkte jeweils einen Abteilungschef bestimmt, welcher die jeweilige Abteilung im Auge behält und sich auch hauptsächlich darum kümmert. In diesem Fall sah die Aufteilung wie folgt aus:

\begin{tabular}{lll}
	\toprule
	\textbf{Name} & \textbf{Abteilung 1} & \textbf{Abteilung 2} \\
	\midrule
	Maurice Manz & mechanische Konstruktion & elektronische Steuerung \\
	Ferdinand Leybach & Softwareentwicklung & Bauteilbeschaffung \\
	Yannick Schmitt & Test und Qualifikation & \\
	Simon Straub & Projektmanagement & \\
	\bottomrule
\end{tabular}




 Ziel ist es, die technische Umsetzung nachvollziehbar darzustellen und die zugrunde liegenden Überlegungen transparent zu machen. Es wird durchaus die Absicht verfolgt nach der Entwicklung des ersten Prototypen auch mehrere dieser Bierkühler zu bauen, sodass einige Leute von unserer Idee profitieren können. Dabei muss natürlich auf auf den Preis geschaut werden, sodass es sich nicht um zu teure Lösungen für kühles Bier handelt.